%==========================================================
% BACKMATTER — Volume I: Zero to Bit
%==========================================================
%
% Structure (per-volume):
%   \appendix  — lettered chapters specific to this volume
%   \backmatter — unnumbered closing material
%
% Shared backmatter (bibliography, index) is handled in main.tex
% after the @BACKMATTER@ token.
%==========================================================

\appendix

%----------------------------------------------------------
\chapter{Standard C and POSIX Headers}
\label{app:headers}
\addcontentsline{toc}{chapter}{Appendix A\quad Standard C and POSIX Headers}
\chaptermark{Standard Headers}
%----------------------------------------------------------

\begin{chapterintro}
The topics in this volume correspond to a well-defined set of standard C
and POSIX headers. This appendix maps those topics to the headers where
the real definitions live, so that moving from the text to actual code is
straightforward. Entries are grouped to match the order of the chapters.
\end{chapterintro}

%---------------------------
\section{Integer Types and Limits}
%---------------------------

\begin{description}[leftmargin=0pt, itemsep=10pt]

\item[\texttt{<stdint.h>} \textnormal{(C99)}]
Exact-width integer types: \texttt{int8\_t}, \texttt{uint8\_t},
\texttt{int16\_t}, \texttt{uint16\_t}, \texttt{int32\_t}, \texttt{uint32\_t},
\texttt{int64\_t}, \texttt{uint64\_t}.
Also defines minimum-width types (\texttt{int\_least8\_t}, \ldots),
fastest minimum-width types (\texttt{int\_fast8\_t}, \ldots),
and pointer-width integers (\texttt{intptr\_t}, \texttt{uintptr\_t}).
Use these names rather than \texttt{int} or \texttt{long} whenever the
exact bit-width is part of the contract.

\item[\texttt{<inttypes.h>} \textnormal{(C99)}]
Portable format macros for \texttt{printf} and \texttt{scanf} with
\texttt{<stdint.h>} types: \texttt{PRId32}, \texttt{PRIu64},
\texttt{SCNi32}, and so on. These expand to the correct format specifier
for each type on each platform.

\item[\texttt{<limits.h>}]
Range limits for the built-in types: \texttt{CHAR\_BIT} (almost always 8),
\texttt{SCHAR\_MIN}, \texttt{SCHAR\_MAX}, \texttt{INT\_MIN}, \texttt{INT\_MAX},
\texttt{UINT\_MAX}, \texttt{LLONG\_MIN}, \texttt{ULLONG\_MAX}, and the rest.
Never hard-code these values; always read them from this header.

\item[\texttt{<stddef.h>}]
\texttt{size\_t} (the unsigned result type of \texttt{sizeof}),
\texttt{ptrdiff\_t} (the signed difference of two pointers), \texttt{NULL},
and the \texttt{offsetof} macro.

\end{description}

%---------------------------
\section{Floating-Point}
%---------------------------

\begin{description}[leftmargin=0pt, itemsep=10pt]

\item[\texttt{<float.h>}]
Implementation characteristics of the floating-point types.
\texttt{FLT\_RADIX} is 2 on any IEEE~754 system.
\texttt{FLT\_MANT\_DIG} and \texttt{DBL\_MANT\_DIG} give the number of
significand bits (24 and 53 for single and double precision).
\texttt{FLT\_EPSILON} and \texttt{DBL\_EPSILON} are machine epsilon.
\texttt{FLT\_DIG} and \texttt{DBL\_DIG} are the number of significant
decimal digits.

\item[\texttt{<math.h>}]
The macros \texttt{INFINITY} and \texttt{NAN}, and the classification
functions \texttt{isnan()}, \texttt{isinf()}, \texttt{isfinite()},
\texttt{isnormal()}, \texttt{fpclassify()}, and \texttt{signbit()}.
These are the standard interface for the special values covered in
Chapters~10--11.

\item[\texttt{<fenv.h>} \textnormal{(C99)}]
Floating-point environment control. \texttt{fegetround()} and
\texttt{fesetround()} read and write the active rounding mode
(\texttt{FE\_TONEAREST}, \texttt{FE\_TOWARDZERO}, \texttt{FE\_UPWARD},
\texttt{FE\_DOWNWARD}). \texttt{fetestexcept()} and \texttt{feclearexcept()}
manage the exception flags: \texttt{FE\_OVERFLOW}, \texttt{FE\_UNDERFLOW},
\texttt{FE\_INEXACT}, \texttt{FE\_DIVBYZERO}, \texttt{FE\_INVALID}.

\end{description}

%---------------------------
\section{Characters and Text}
%---------------------------

\begin{description}[leftmargin=0pt, itemsep=10pt]

\item[\texttt{<ctype.h>}]
Character classification (\texttt{isalpha()}, \texttt{isdigit()},
\texttt{isspace()}, \texttt{isprint()}, \ldots) and case conversion
(\texttt{toupper()}, \texttt{tolower()}) for single-byte characters.
All functions take and return \texttt{unsigned char} values, not plain
\texttt{char}. Behaviour is locale-dependent.

\item[\texttt{<wchar.h>} \textnormal{(C95)}]
Wide-character strings: \texttt{wcslen()}, \texttt{wcscpy()},
\texttt{wprintf()}, \texttt{fgetwc()}. Multibyte-to-wide-character
conversion: \texttt{mbstowcs()}, \texttt{wcstombs()}, \texttt{mbtowc()},
\texttt{mbrtowc()}.

\item[\texttt{<uchar.h>} \textnormal{(C11)}]
\texttt{char16\_t} and \texttt{char32\_t} for UTF-16 and UTF-32 code units.
Restartable conversion functions: \texttt{mbrtoc16()}, \texttt{c16rtomb()},
\texttt{mbrtoc32()}, \texttt{c32rtomb()}.

\end{description}

%---------------------------
\section{Byte Order}
%---------------------------

Standard C does not specify the byte order of multi-byte integers, and
before C23 there is no standard header for host-byte-order conversion.
The following platform-specific interfaces cover most cases:

\begin{description}[leftmargin=0pt, itemsep=10pt]

\item[\texttt{<endian.h>} \textnormal{(Linux / glibc)}]
\texttt{htobe16()}, \texttt{htobe32()}, \texttt{htobe64()} convert from
host byte order to big-endian; \texttt{be16toh()}, \texttt{be32toh()},
\texttt{be64toh()} reverse the direction. Little-endian counterparts
(\texttt{htole16()}, \texttt{le32toh()}, \ldots) are also provided.

\item[\texttt{<sys/endian.h>} \textnormal{(BSD, macOS)}]
The same functions under a different header name.

\item[\texttt{<arpa/inet.h>} \textnormal{(POSIX)}]
The original network-byte-order functions: \texttt{htonl()}, \texttt{htons()},
\texttt{ntohl()}, \texttt{ntohs()}. Available on every POSIX system and
still the most portable choice.

\end{description}

%---------------------------
\section{Bit Operations}
%---------------------------

\begin{description}[leftmargin=0pt, itemsep=10pt]

\item[GCC and Clang built-ins]
Standard C before C23 provides no bit-manipulation functions. The
following compiler built-ins are available on GCC, Clang, and compatible
compilers. Each has an \texttt{l} (long) and \texttt{ll} (long long)
variant:
\begin{itemize}[noitemsep]
  \item \texttt{\_\_builtin\_popcount(x)} --- number of set bits
  \item \texttt{\_\_builtin\_clz(x)} --- leading zero count (undefined if \texttt{x} is zero)
  \item \texttt{\_\_builtin\_ctz(x)} --- trailing zero count (undefined if \texttt{x} is zero)
  \item \texttt{\_\_builtin\_parity(x)} --- XOR of all bits (1 if odd number of set bits)
\end{itemize}

\item[\texttt{<stdbit.h>} \textnormal{(C23)}]
The first standard bit-manipulation header.
Provides \texttt{stdc\_popcount()}, \texttt{stdc\_leading\_zeros()},
\texttt{stdc\_trailing\_zeros()}, \texttt{stdc\_bit\_width()},
\texttt{stdc\_has\_single\_bit()}, \texttt{stdc\_bit\_floor()},
and \texttt{stdc\_bit\_ceil()}. Prefer these over compiler built-ins
when your toolchain supports C23.

\end{description}

%---------------------------
\section{Memory and Alignment}
%---------------------------

\begin{description}[leftmargin=0pt, itemsep=10pt]

\item[\texttt{<string.h>}]
Primitive operations on raw byte arrays: \texttt{memcpy()},
\texttt{memmove()}, \texttt{memset()}, \texttt{memcmp()}, \texttt{memchr()}.

\item[\texttt{<stdlib.h>}]
Dynamic allocation: \texttt{malloc()}, \texttt{calloc()},
\texttt{realloc()}, \texttt{free()}.
C11 adds \texttt{aligned\_alloc(alignment, size)} for memory with a
guaranteed alignment requirement.

\item[\texttt{<stdalign.h>} \textnormal{(C11)}]
\texttt{alignof(T)} returns the alignment requirement of type \texttt{T}.
\texttt{alignas(N)} specifies the alignment for a variable or type member.
Both became keywords in C23; the header provides them as macros for C11
and C17.

\end{description}

%---------------------------
\section{Error Detection and I/O}
%---------------------------

\begin{description}[leftmargin=0pt, itemsep=10pt]

\item[\texttt{<stdio.h>}]
The format specifiers for printing integers and floating-point values.
When working with \texttt{<stdint.h>} types, use the macros from
\texttt{<inttypes.h>} rather than hard-coded specifiers such as
\texttt{\%d} or \texttt{\%lu}.

\item[\texttt{<errno.h>}]
Error codes for library functions: \texttt{ERANGE} (result out of
representable range), \texttt{EDOM} (domain error in a math function),
and others. The global variable \texttt{errno} is set by library
functions on failure.

\end{description}

%==========================================================
% End of Volume I closing
%==========================================================
\backmatter
\pagestyle{backmatter}

\chapter*{End of Volume\,I}
\addcontentsline{toc}{chapter}{End of Volume I}
\chaptermark{End of Volume I}

Volume~I has followed one thread: how information is encoded in digital
systems, starting at a voltage level and ending at compressed binary
streams. Integers, floating-point numbers, Unicode text, byte order,
aligned memory, pointers, error-detection codes, serialization formats ---
these are not separate topics. They are successive layers of agreed-upon
meaning stacked on top of the same physical signal.

Volume~II picks up where this one ends. The encodings you have read about
here exist in physical hardware --- transistors, logic gates, memory cells,
cache hierarchies --- that imposes its own constraints on how those encodings
are stored, moved, and transformed. Understanding those constraints is what
turns a working program into a fast one.

\medskip
\noindent\textit{If any chapter in this volume felt unclear, come back to it
after Volume~II. The hardware context makes many representational choices
obvious in retrospect.}

\clearpage